% Part: second-order-logic
% Chapter: syntax-and-semantics
% Section: introduction

\documentclass[../../../include/open-logic-section]{subfiles}

\begin{document}

\olfileid[hi]{sol}{syn}{int}

\olsection{परिचय}

प्रथम-कोटि तर्कशास्त्र में हम किसी दी हुई भाषा के अतार्किक चिह्नों, अर्थात्
उसके !!{constant}, !!{function} और !!{predicate} चिह्नों को तार्किक चिह्नों
के साथ जोड़कर प्रथम-कोटि !!{structure} के विषय में बातें व्यक्त करते हैं। यह
कार्य संतुष्टि की धारणा से किया जाता है, जो !!a{structure}~$\Struct{M}$,
चर-नियतन~$s$ और !!a{formula}~$!A$ को परस्पर संबद्ध करती है:
$\Sat{M}{!A}[s]$ तभी लागू होता है जब $!A$ द्वारा व्यक्त वह बात सत्य हो जिसमें
उसके !!{constant}, !!{function} और !!{predicate} चिह्नों की व्याख्या
$\Struct{M}$ के अनुसार तथा उसके मुक्त चरों की व्याख्या~$s$ के अनुसार की गई
हो। !!{identity}~$\eq$ की व्याख्या $\Sat{M}{!A}[s]$ की परिभाषा में ही
अंतर्निहित है; यही बात~$\lforall$ और $\lexists$ की व्याख्या पर भी लागू होती
है। पहले की व्याख्या संरचना के !!{domain}~$\Domain{M}$ पर सर्वदा तादात्म्य
संबंध के रूप में होती है, और परिमाणक सर्वदा पूरे !!{domain} पर विचरते हैं।
किंतु महत्त्वपूर्ण बात यह है कि परिमाणीकरण केवल !!{domain} के अवयवों पर
स्वीकृत है; अतः परिमाणक के बाद केवल वस्तु-!!{variable} ही आ सकते हैं।

द्वितीय-कोटि तर्कशास्त्र में भाषा और संतुष्टि की परिभाषा, दोनों का विस्तार
इस प्रकार किया जाता है कि उनमें मुक्त और आबद्ध फलन-चर तथा विधेय-चर और उन पर
परिमाणीकरण सम्मिलित हो जाएँ। इन चरों का !!{function} और !!{predicate} से वही
संबंध है जो वस्तु-चरों का !!{constant} से है। द्वितीय-कोटि तर्कशास्त्र के पदों
और !!{formula} के निर्माण में ये वही भूमिका निभाते हैं, और इन पर परिमाणीकरण
भी इसी प्रकार किया जाता है। \emph{मानक} अर्थविज्ञान में द्वितीय-कोटि
परिमाणक उचित प्रकार की सभी संभव वस्तुओं पर विचरते हैं (फलन-चरों के लिए
$\Domain{M}$ से~$\Domain{M}$ में जाने वाले $n$-स्थानी फलन और विधेय-चरों के
लिए $n$-स्थानी संबंध)। उदाहरणार्थ,
$\lforall[\Obj{v_0}][(\Obj{P^1_0}(\Obj{v_0}) \lor \lnot
  \Obj{P^1_0}(\Obj{v_0}))]$ प्रथम- और द्वितीय-कोटि, दोनों तर्कशास्त्रों में एक सूत्र है; किंतु दूसरे में हम
इन पर भी विचार कर सकते हैं:
$\lforall[\Obj{V^1_0}][\lforall[\Obj{v_0}][(\Obj{V^1_0}(\Obj{v_0}) \lor
    \lnot \Obj{V^1_0}(\Obj{v_0}))]]$ और
$\lexists[\Obj{V^1_0}][\lforall[\Obj{v_0}][(\Obj{V^1_0}(\Obj{v_0}) \lor
    \lnot \Obj{V^1_0}(\Obj{v_0}))]]$। इनमें कोई मुक्त चर नहीं है, इसलिए ये
द्वितीय-कोटि तर्कशास्त्र के !!{sentence} हैं। यहाँ $\Obj{V^1_0}$ एक
द्वितीय-कोटि $1$-स्थानी विधेय-चर है। $\Obj{V^1_0}$ की स्वीकार्य व्याख्याएँ
वही हैं जिन्हें हम $\Obj{P^1_0}$ जैसे $1$-स्थानी !!{predicate} को दे सकते
हैं, अर्थात्~$\Domain{M}$ के उपसमुच्चय। तब उन पर परिमाणीकरण का अर्थ यह कहना
है कि
$\lforall[\Obj{v_0}][(\Obj{V^1_0}(\Obj v_0) \lor \lnot
  \Obj{V^1_0}(v_0))]$ उन सभी विधियों के लिए लागू होता है जिनसे
$\Obj{V^1_0}$ का मान~$\Domain{M}$ का कोई उपसमुच्चय नियत किया जा सकता है,
अथवा कम-से-कम एक ऐसी विधि के लिए लागू होता है। चूँकि प्रत्येक समुच्चय में
दी हुई वस्तु या तो होती है या नहीं होती, इसलिए दोनों किसी भी !!{structure}
में सत्य हैं।
\end{document}
